\chapter{Suspension Objects}
\label{ch:suspension}

A \emph{suspension object} is the lightest synchronisation primitive Ada offers: a
single Boolean flag with one special power --- a task can \emph{block} on it until
the flag is set, consuming the signal as it wakes. It is the idiomatic way to hand
an event from an interrupt handler to a normal task (``the edge fired, go look''),
and it costs almost nothing. The tasking chapter introduced it in passing while
surveying ways to wait for an event (\S\ref{sec:waiting}); this chapter is the
focused account --- the exact semantics, the rules that bite, and when to reach for
one instead of a protected entry.

Suspension objects live in \texttt{Ada.Synchronous\_Task\_Control} (ARM~D.10) and
are available under \emph{every} profile --- \texttt{light-tasking},
\texttt{embedded}, and \texttt{full} alike --- because they are a core Ravenscar
primitive.

\section{The interface}

The whole package is four operations over one opaque, limited type:

\begin{lstlisting}
package Ada.Synchronous_Task_Control with Preelaborate is
   type Suspension_Object is limited private;     --  one bit; starts False

   procedure Set_True  (S : in out Suspension_Object);   --  raise the flag
   procedure Set_False (S : in out Suspension_Object);   --  lower it
   function  Current_State (S : Suspension_Object) return Boolean;
   procedure Suspend_Until_True (S : in out Suspension_Object);  --  block + consume
end Ada.Synchronous_Task_Control;
\end{lstlisting}

\texttt{Suspend\_Until\_True} is the one with teeth. If the flag is already
\texttt{True} it returns at once, leaving it \texttt{False}; if it is \texttt{False}
the calling task \emph{blocks} --- descheduled, burning no CPU --- until some other
task (typically an interrupt handler) calls \texttt{Set\_True}, at which point it
wakes and, in the same indivisible step, clears the flag back to \texttt{False}.

There is no magic inside. On the bare runtime a \texttt{Suspension\_Object} is just
a protected object with a one-line barrier, and the whole of D.10 is this:

\begin{lstlisting}
protected body Suspension_Object is
   procedure Set_True  is begin Open := True;  end Set_True;
   procedure Set_False is begin Open := False; end Set_False;

   entry Wait when Open is        --  Suspend_Until_True calls this entry
   begin
      Open := False;              --  consume the signal as we pass the barrier
   end Wait;
end Suspension_Object;
\end{lstlisting}

Seeing the implementation explains every rule that follows.

\section{The signalling pattern}

The canonical use is interrupt-to-task hand-off. A protected interrupt handler does
the absolute minimum --- set the flag --- and a worker task does the real work after
\texttt{Suspend\_Until\_True} wakes it. The handler stays tiny (it runs at a high
interrupt ceiling, \S\ref{ch:interrupts}); the work runs at ordinary task priority.

\begin{lstlisting}
with Ada.Synchronous_Task_Control;  use Ada.Synchronous_Task_Control;
with ESP32S3.GPIO;
with ESP32S3.GPIO.Interrupts;

Edge_Seen : Suspension_Object;       --  library level

procedure On_Edge is                 --  runs in interrupt context: keep it tiny
begin
   Set_True (Edge_Seen);             --  "an edge happened"
end On_Edge;

task Worker;
task body Worker is
begin
   ESP32S3.GPIO.Configure (0, ESP32S3.GPIO.Input, Pull => ESP32S3.GPIO.Pull_Up);
   ESP32S3.GPIO.Interrupts.Enable
     (Pin => 0, On => ESP32S3.GPIO.Interrupts.Falling_Edge, Action => On_Edge'Access);
   loop
      Suspend_Until_True (Edge_Seen);   --  sleep here, zero CPU, until an edge
      --  ... handle the edge ...
   end loop;
end Worker;
\end{lstlisting}

This is exactly how the HAL's own interrupt children (GPIO, and the
\texttt{*-interrupts} packages for the I2C devices) deliver their events.

\section{Level-triggered, so no lost wake-ups}

The flag is \emph{sticky}. A \texttt{Set\_True} that happens \emph{before} the task
reaches \texttt{Suspend\_Until\_True} is not lost: the task finds the flag already
\texttt{True} and returns immediately, consuming it. That closes the classic
lost-wake-up race --- the one a raw ``wake'' signal or a bare condition variable
suffers when the signal arrives a moment before the wait. With a suspension object
the order does not matter.

The flip side is that it is a \emph{binary flag, not a counter}. Several
\texttt{Set\_True} calls before a single \texttt{Suspend\_Until\_True} collapse into
\emph{one} wake-up. A suspension object therefore means ``at least one event has
happened since you last looked,'' never ``exactly $N$ events.'' For a stream where
every event must be accounted for, that coalescing is a bug waiting to happen --- see
the protected-entry upgrade below.

\section{One waiter, and only self-suspension}

Two rules fall straight out of that \texttt{entry Wait}:

\begin{itemize}
\item \textbf{A task may only suspend \emph{itself}.} \texttt{Suspend\_Until\_True}
  blocks the caller; there is no operation to suspend another task. Any task (or
  handler) may \emph{resume} a blocked one with \texttt{Set\_True}.
\item \textbf{At most one task may wait on a given object at a time.} A second task
  that calls \texttt{Suspend\_Until\_True} while another is already blocked raises
  \texttt{Program\_Error}. This is not an arbitrary limit: the object is a protected
  object with a single entry, and the Jorvik/Ravenscar restriction
  \texttt{Max\_Entry\_Queue\_Length =>\ 1} caps that entry's queue at one caller.
  One suspension object serves one waiter.
\end{itemize}

If several tasks must wake on the same event, give each its own suspension object
and signal them in turn, or use a protected object whose barrier they all share.

\section{Suspension object or protected entry?}

A suspension object carries \emph{no data} --- it only says ``something happened.''
When you need to pass a value across the hand-off, or to count events rather than
coalesce them, the natural upgrade is a protected object with an entry: the same
block-until-ready behaviour, but the barrier can be any condition and the entry can
deliver a payload. The coalescing problem above, fixed, is just such an entry:

\begin{lstlisting}
protected Edge_Box is
   procedure Note;                       --  called from the ISR callback
   entry     Wait (Count : out Natural); --  block until >= 1, return how many
private
   N : Natural := 0;
end Edge_Box;

protected body Edge_Box is
   procedure Note is begin N := N + 1; end Note;       --  never loses a count
   entry Wait (Count : out Natural) when N > 0 is
   begin
      Count := N;  N := 0;               --  consume all of them, atomically
   end Wait;
end Edge_Box;
\end{lstlisting}

So: reach for a \textbf{suspension object} when the signal is simply ``it happened,
go and look,'' and the latest state is what matters (a DMA completion, a button
edge, ``data is ready in the shared buffer''). Reach for a \textbf{protected entry}
when you must not lose counts, must transfer a value atomically, or must block on a
richer condition than one bit. The protected-entry half of this trade is developed
further in the tasking chapter's waiting patterns (\S\ref{sec:waiting}).

\section{No timeout of its own}

\texttt{Suspend\_Until\_True} blocks until the flag is set, full stop --- there is no
``wait, but give up after $T$'' form. When you need a bounded wait, the approach
depends on the profile, and the tasking chapter works all three through in detail
(\S\ref{sec:waiting}): under Jorvik a protected entry barriered on the event
\emph{or} a deadline (an \texttt{Ada.Real\_Time.Timing\_Events} handler); under
\texttt{full} a \texttt{select \ldots\ Suspend\_Until\_True; or delay until \ldots}
asynchronous transfer of control. For the simple ``poll the hardware but give up''
case that needs no interrupt at all, \texttt{ESP32S3.Wait.Until\_True} is a
lock-free helper that works under every profile.

\section{Cost and when it shines}

A suspension object is about the cheapest blocking primitive there is: one protected
object, a one-line barrier, no payload to copy, no ceiling negotiation beyond the
protected action itself. \texttt{Set\_True} from an interrupt handler is a couple of
instructions, and the woken task was genuinely asleep --- no spinning, no polling,
no wasted power --- until the hardware had something to say. For the overwhelmingly
common embedded shape ``a short ISR notices an event; a normal task does the work,''
it is the right tool, and the one the HAL reaches for first.
